\section{Metrics, formulations, and the finite-parameter boundary}
\label{sec:formulation-scope}

\subsection{Fixed coordinate changes and gap-dependent preconditioning}

All preceding spectral rates use a fixed external metric on the
feasible affine hull. If $R:\R^d\to\R^d$ is a fixed invertible
change of tangent coordinates, the Hessian matrix becomes $R^\top HR$.
Using ordinary Euclidean matrix condition numbers in the respective
coordinate systems gives
\begin{equation}
 \frac{\kappa(H)}{\kappa(R)^2}
 \leq\kappa(R^\top HR)\leq\kappa(R)^2\kappa(H).
 \label{eq:coordinate-condition-comparison}
\end{equation}
The upper bound follows by bounding the extreme Rayleigh quotients
using the extreme singular values of $R$; the lower bound follows by
applying the upper bound to $R^{-1}$. Thus fixed coordinate changes
preserve accuracy exponents, with possibly large constant factors.
If the original external metric is transported instead, its Gram matrix
is $R^\top R$, and the generalized spectrum of
$(R^\top HR,R^\top R)$ is exactly the spectrum of $H$.
A gap-dependent transformation is different: the mathematical choice
$R(g)=\widehat H_F(g)^{-1/2}$ produces the identity as the Euclidean
matrix for the transformed linear solve. Whether such a
preconditioner is useful depends on its construction, application,
and representation costs. Neither the geometric theorem nor
\eqref{eq:coordinate-condition-comparison} bounds those costs.

\subsection{The reduced Hessian is not the Schur or augmented matrix}

To distinguish the matrices, let $B=\nabla^2F(x)\succ0$ be an
\emph{ambient} barrier Hessian for equality constraints $Ax=b$.
The reduced matrix is $W^\top BW$ for an orthonormal basis of
$\ker A$; the primal Newton Schur matrix is $AB^{-1}A^\top$.
The latter expression must use $B$, not the already reduced
operator. Under an invertible ambient coordinate change $z=Gx$,
transport the barrier as $\widetilde F(z)=F(G^{-1}z)$ and put
$\widetilde A=AG^{-1}$. Then
\begin{equation}
 \widetilde B=G^{-\top}BG^{-1},\qquad
 \widetilde A\widetilde B^{-1}\widetilde A^\top
 =AB^{-1}A^\top.
 \label{eq:schur-invariance}
\end{equation}
In contrast, the augmented Newton matrix transforms by congruence:
\begin{equation}
 \begin{pmatrix}\widetilde B&\widetilde A^\top\\
                 \widetilde A&0\end{pmatrix}
 =\begin{pmatrix}G^{-\top}&0\\0&I\end{pmatrix}
   \begin{pmatrix}B&A^\top\\A&0\end{pmatrix}
   \begin{pmatrix}G^{-1}&0\\0&I\end{pmatrix}.
 \label{eq:kkt-congruence}
\end{equation}
Congruence is not spectral equality. In particular, a bound for the
reduced primal Hessian cannot be read as the same bound for every
Schur, scaled primal--dual, or augmented formulation. Classical
conditioning analyses such as \citet{pena2002,renegar1996} must be
compared with the actual matrix they study.

For an explicit illustration of \eqref{eq:schur-invariance}, consider
the Lorentz cone and its barrier
$F(t,z)=-\log(t^2-\norm z^2)$. For $0<\lambda<1$, let $G_\lambda$
act on the first two coordinates as
\[
 \frac12\begin{pmatrix}
 \lambda+\lambda^{-1}&\lambda-\lambda^{-1}\\
 \lambda-\lambda^{-1}&\lambda+\lambda^{-1}
 \end{pmatrix}
\]
and as the identity on the other coordinates. It preserves
$t^2-\norm z^2$ and the positive sheet, so
$F(G_\lambda x)=F(x)$. At $e_0=(1,0,\ldots,0)$ the ambient Hessian
is $2I$, while at $G_\lambda e_0$ it is $2G_\lambda^{-2}$, with
condition number $\lambda^{-4}$. Nevertheless, transforming $A$
as above leaves the Schur matrix exactly unchanged. This is a
coordinate illustration on the cone, not an application of the
compact-feasible-set theorem to that unbounded cone. A varying
$\lambda$ can also change coefficient magnitudes and encoding
normalizations, even though the transformation mixes only two
coordinates.

The intrinsic length of a differentiable curve,
$\int\sqrt{\dot x^\top B(x)\dot x}\,dt$, is preserved by the same
transported metric. Its geometric role is studied by
\citet{nesterovtodd2002}. This integral and a pointwise external
condition number describe different quantities. Sensitivity to the
objective also requires specifying an absolute scale. At fixed
$\mu$, implicit differentiation of exact centrality gives
\begin{equation}
 D_{c_{\V}}x_F(\mu)=-\mu^{-1}H_F(x_F(\mu))^{-1}.
 \label{eq:objective-sensitivity}
\end{equation}
Thus a lower Hessian eigenvalue bound does control this absolute
objective-to-center derivative; the condition number alone does not.
A derivative with respect to a chosen family parameter additionally
includes the derivative of the data. For example, choose a short line segment
$z_0+[-a,a]v$ contained in a fixed barrier domain and set
$z_k(t)=z_0+a\sin(kt)v$. Continuity and positive definiteness give
uniform spectral bounds for the Hessian on that segment, while
$\norm{z_k'(0)}=ak\norm v$ is unbounded. Taking
$c_k(t)=-\nabla F(z_k(t))$ makes each point an exact center for
$\mu=1$ with that varying objective. The data derivatives vary too,
consistently with \eqref{eq:objective-sensitivity}. The example
establishes a separation between pointwise spectral bounds and variation
with respect to an unrestricted family parameter. It does not show
unbounded sensitivity per unit objective perturbation when the Hessian
is uniformly bounded below. More generally, differentiating an auxiliary
representation of a central point introduces derivatives of that
representation's data maps. Comparing its sensitivity with Newton
conditioning requires bounds on those maps and on the Jacobian of its
defining equations; pointwise centrality alone does not supply that
comparison.

\subsection{Why a finite barrier parameter is essential}

The gradient inequality \eqref{eq:barrier-gradient} is essential to
the containment theorem. Differential self-concordance
\eqref{eq:sc} and boundary divergence alone do not suffice.
On the same rectangle as \cref{prop:oscillatory-barrier}, for fixed
$\eta>0$ consider
\begin{equation}
 F_\eta(x,y)=-\log(1-x^2)-\log y-\log(1-y)+\frac{\eta}{y}.
 \label{eq:infinite-parameter-example}
\end{equation}
This function is convex, diverges at the boundary, and is
differentially self-concordant. To check the last assertion, for
$h(y)=-\log y+\eta/y$ write $t=\eta/y\geq0$. Then
\[
 h''(y)=y^{-2}(1+2t),\qquad
 |h'''(y)|=2y^{-3}(1+3t),\qquad
 (1+3t)^2\leq(1+2t)^3,
\]
where the difference of the right and left sides is $3t^2+8t^3$.
Adding the other affine logarithms preserves differential
self-concordance, by summing their third-derivative bounds and using
$\sum a_i^{3/2}\leq(\sum a_i)^{3/2}$ for $a_i\geq0$.

For objective $y$, its central tail has $x=0$ and gap $g=y$.
Indeed $\partial_xF_\eta(0,g)=0$ and
$\partial_yF_\eta(0,g)=-1/g+1/(1-g)-\eta/g^2<0$ for small $g$,
so $\mu(g)=-1/\partial_yF_\eta(0,g)>0$ makes $(0,g)$ the unique
center by strict convexity.
The diagonal Hessian has entries
$2$ and $g^{-2}+(1-g)^{-2}+2\eta g^{-3}$, so
\begin{equation}
 \widehat\kappa_{F_\eta}(g)\sim\eta g^{-3},\qquad
 \frac{|\partial_yF_\eta(0,g)|^2}
      {\partial_{yy}F_\eta(0,g)}\sim\frac{\eta}{2g}\longrightarrow\infty.
 \label{eq:infinite-parameter-rate}
\end{equation}
Hence no finite $\nu$ satisfies the barrier-gradient inequality.
This does not contradict the $\Theta(g^{-2})$ finite-parameter law
for a positive-dimensional optimal face. The same distinction
between differential self-concordance and a finite global gradient
parameter is explicitly noted for inverse-power perturbations in
\citet[Remark 11.1]{monteirodasilva2026}; no result of that preprint
is needed for the calculation here.
